%!TEX program = lualatex
% 样稿：与正文（latex/book.tex）共用前导、前件与章节文件，
% 唯一区别是样稿标题页与样稿说明。
\documentclass[11pt,a4paper,fontset=fandol]{ctexbook}

%!TEX program = lualatex
% 共享前导：latex/book.tex（正文）与 sample/sample.tex（样稿）共用。

\usepackage{amsmath,amssymb,amsthm,bm}
\usepackage{etoolbox}
\usepackage[left=3.2cm, right=3cm, top=3cm, bottom=3cm]{geometry}
\usepackage{booktabs,longtable,array}
\usepackage{calc}
\newcounter{none} % pandoc longtable without caption
\usepackage{graphicx,float}
\usepackage{enumerate,enumitem}
\usepackage[font=small]{caption}
\usepackage{xcolor}
\usepackage{tikz}
\usetikzlibrary{matrix,arrows,positioning,fit,calc}
\usepackage{listings}
\usepackage{fancyvrb}
\usepackage{newunicodechar}
\usepackage[unicode=true,pdfencoding=auto,psdextra,
  colorlinks=true,linkcolor=blue,urlcolor=blue]{hyperref}
\makeatletter
\g@addto@macro{\UrlBreaks}{\UrlOrds}
\def\input@path{{../latex/}{../latex/chapters/}{chapters/}}
\makeatother

\setlength{\parindent}{2em}

\providecommand{\tightlist}{%
  \setlength{\itemsep}{0pt}\setlength{\parskip}{0pt}}
\newcommand{\passthrough}[1]{#1}

\lstset{
  basicstyle=\ttfamily\small,
  breaklines=true,
  columns=flexible,
  keepspaces=true,
  showstringspaces=false,
  aboveskip=6pt,
  belowskip=6pt,
  frame=single,
  framesep=4pt,
  rulecolor=\color{black!30},
  keywordstyle=\color{blue!45!black}\bfseries,
  commentstyle=\color{black!45}\itshape,
  stringstyle=\color{red!55!black},
  literate={→}{{$\to$}}1 {←}{{$\leftarrow$}}1 {⇒}{{$\Rightarrow$}}1
           {≠}{{$\neq$}}1 {≥}{{$\geq$}}1 {≤}{{$\leq$}}1 {∈}{{$\in$}}1,
}

% Markers that have no glyph in the CJK fonts.
% 🦂 ◇ ✗ ⚠ render as color emoji via the emoji package (LuaLaTeX; Apple Color
% Emoji auto-detected).
\usepackage{emoji}
\newunicodechar{🦂}{\emoji{scorpion}}
\newunicodechar{◇}{\emoji{large-blue-diamond}}
\newunicodechar{✗}{\emoji{cross-mark}}
\newunicodechar{⚠}{\emoji{warning}}
\newunicodechar{→}{\textrightarrow}
\newunicodechar{←}{\textleftarrow}
\newunicodechar{⇒}{\ensuremath{\Rightarrow}}
\newunicodechar{≠}{\ensuremath{\neq}}
\newunicodechar{≥}{\ensuremath{\geq}}
\newunicodechar{≤}{\ensuremath{\leq}}
\newunicodechar{∈}{\ensuremath{\in}}
\newunicodechar{∅}{\ensuremath{\varnothing}}
\newunicodechar{μ}{\ensuremath{\mu}}
\newunicodechar{λ}{\ensuremath{\lambda}}


\begin{document}

\begin{titlepage}
  \centering
  \vspace*{3.2cm}
  {\Large\bfseries 边学 Haskell 边学范畴论\par}
  \vspace{0.4cm}
  {\large\bfseries 从可运行的代码反推数学\par}
  \vspace{1.2cm}
  {\large 宋丛威\par}
  \vspace{2.2cm}
  {\large 样稿\par}
  \vspace{0.4cm}
  {\normalsize 内容简介 $\cdot$ 序 $\cdot$ 目录 $\cdot$ 第一章 $\cdot$ 第五章（全文）$\cdot$ 其余章节（仅标题）\par}
  \vspace{1.2cm}
  {\small 供编辑与评审审阅\par}
\end{titlepage}

\frontmatter

\chapter*{样稿说明}
\addcontentsline{toc}{chapter}{样稿说明}

本册为《边学 Haskell 边学范畴论》的样稿，供编辑与评审参考。样稿与正式版正文的唯一区别
是本说明与样稿标题页；正文取自正式版对应章节，未作删改。样稿包含内容简介、序、目录、
第一章与第五章的全文；其余章节在目录中列出标题，正文从略。

\chapter*{内容简介}
\addcontentsline{toc}{chapter}{内容简介}


本书面向会 Python、写过生产代码、却被 Haskell 类型系统与范畴论术语卡住的程序员。
它不把数学当作语言的背景知识，而是从能跑的代码倒推数学：每一章先给出一个写起来难看
的程序，再由它逼出一个概念。函子、应用函子、单子因此不再是名词，而是三个具体问题的
答案：给不同容器写同一套变换、把二元函数放进语境、三种效应共享一套复合规则。

全书分五个部分。语言部分用同一个求值器演示 Haskell 与其他语言的差异，并引入代数数据
类型与类型类；范畴与函子部分给出范畴的定义，核对 Hask 满足公理的范围，引入函子与自然
变换；应用函子部分把 Applicative 对应到数学上的 lax monoidal 函子，并巡礼六个实例；单子
部分处理 join 与 Kleisli 范畴；工程部分覆盖输入输出、包与测试、并发和一个完整实战。

本书与同类书的差别在于可验证：正文的每一段代码都来自可编译的工程，由测试与 CI 逐条
核对；数学与代码不一致处（如 $\bot$ 使 Hask 不满足单位律）不回避，明说结论成立的范围。


\chapter*{序}
\addcontentsline{toc}{chapter}{序}

\input{chapters/preface}

\tableofcontents

\mainmatter
\setcounter{secnumdepth}{3}
\part{语法与特性}

\input{chapters/ch01}

\input{chapters/ch02}

\input{chapters/ch03}

\input{chapters/ch04}

\part{范畴论}

\input{chapters/ch05}

\input{chapters/ch06}

\input{chapters/ch07}

\input{chapters/ch08}

\input{chapters/ch09}

\input{chapters/ch10}

\part{工程}

\input{chapters/ch11}

\input{chapters/ch12}

\input{chapters/ch13}

\part{实战}

\input{chapters/ch14}

\input{chapters/ch15}

\input{chapters/ch16}

\input{chapters/ch17}

\part{进阶}

\input{chapters/ch18}

\input{chapters/ch19}

\input{chapters/ch20}

\input{chapters/ch21}


\end{document}
